\printedpage{39}
% head: A. Grothendieck : Dérivateurs, Ch. XVI, Édition des Universités de Montpellier II et Paris VII

\msp{page 69}
That $\mathbf{W}_0$ satisfies this set of axioms is checked by AQT \texttt{[âne qui trotte: at a donkey's trot, i.e. by routine verification]}. (\textbf{N.B.} That it satisfies Loc (1 ter, 2, 3 bis) is already a consequence of the fact that $\mathbf{W}_0$ is of the form $W_{\mathbf{D}}$, where $\mathbf{D}$ ($=\mathbf{D}_{\mathbf{W}_0}$ in this case) is a prederivator satisfying Der (3, 4, 5), which is indeed the case for $\mathbf{D}_{\mathbf{W}_0}$, which is a genuine \textit{derivator.}) We must check that every $f$ in Cat factors as $f = pi$, with $i \in \mathbf{W}_0$, $p$ a $\mathbf{W}_0$-fibration. Now if $f : X \longrightarrow Y$, $I = \pi_0(X)$, we take $\bar{X} = \coprod_{\alpha\in I} Y_{f(\alpha)}$, whence a factorization $f = pi$, where $\pi_0(i)$ is an isomorphism, \textit{i.e.} $i \in \mathbf{W}_0$, and where $p : \bar{X} \longrightarrow Y$ is, over the connected component $Y_j$ of $Y$, given by the product $\bar{X}|Y_j \simeq Y_j \times \underbrace{\pi_0(f)^{-1}(j)}_{\text{discrete category}}$ . The morphism $p$ is indeed a $\mathbf{W}_0$-fibration, for thanks to Loc (4), it suffices to check this on each $Y_j$, and there it is deduced from an arrow $I_j \longrightarrow e$ (a $\mathbf{W}_0$-fibration) by base change $Y_j \longrightarrow e$.

On the other hand, that $\mathbf{W}_0$ does \textit{not} even satisfy $\mathrm{D}_{W_\omega}$ has already been seen (I, p. 87),
\msp{page 70}
by realizing the map of the $n$-covering $f : S^1 \longrightarrow S^1$ in Cat by a $W_\omega$-perfect morphism (in fact, by a locally trivial Cat-fibration), a very easy thing\footnote[37]{Another way of seeing it: we \textit{know} (XII) that $\mathrm{W}$ (5, 7 $\mathrm{D}_W$) $\Longrightarrow W \subset \mathbf{W}_\infty$, now we have $\mathbf{W}_\infty \subset \mathbf{W}_0$, and the inclusion is \textit{strict\,!}}.

I think that by now the whole set of different variants of the axiom of $W$-fibrations is more or less understood, except that it is still not clear, without the axiom Loc (8 bis), whether the implications in the diagram
\[
\begin{tikzcd}[row sep=large, column sep=large]
\mathrm{A} \arrow[r, Rightarrow] \arrow[d, Rightarrow] & \mathrm{A}_{W_0} \arrow[d, Rightarrow] \\
\mathrm{B} \arrow[r, Rightarrow] \arrow[d, Rightarrow] & \mathrm{B}_{W_0} \arrow[d, Rightarrow] \\
\mathrm{C} \arrow[r, Rightarrow] \arrow[d, Rightarrow] & \mathrm{C}_{W_0} \arrow[d, Rightarrow] \\
\mathrm{D} \arrow[r, Rightarrow] & \mathrm{D}_{W_0}
\end{tikzcd}
\]
are strict or not. But it seems to me that all the interesting basic localizers encountered so far satisfy Loc (8), and even Loc (1 ter, 2, 3 bis, 4, 5, 6 a, 8, 9) \marginpar{?}(perhaps even Loc (7 F) - I would have to test on the $\mathbf{W}_n$), so the unresolved perplexity seems to me of no great consequence. % [sic] résolu -> résolue, accord avec « la perplexité » ; « ? » imprimé dans la marge gauche à la hauteur de cette ligne

\msp{page 71}

\printedpage{40}
% head: A. Grothendieck : Dérivateurs, Ch. XVI, Édition des Universités de Montpellier II et Paris VII

\section*{5\quad Limits axiom Loc (8), and the accessibility axiom Loc (9)}

It is the axiom:

\textbf{Loc (8)}\quad $W$ is stable under filtering $\varinjlim$.

A useful weakened variant is the following:

\textbf{Loc (8 bis)}\quad $W$ is stable under countable $\varinjlim$.

Explicitly, if we have two filtering direct systems in Cat
\[
\begin{aligned}
\underline{X} &= (X_i)_{i\in I}\\
\underline{Y} &= (Y_i)_{i\in I},
\end{aligned}
\]
and an arrow
\[
\underline{f} = (f_i)_{i\in I} : \underline{X} \longrightarrow \underline{Y},
\]
given by $f_i : X_i \longrightarrow Y_i$, putting $X = \varinjlim X_i$, $Y = \varinjlim Y_i$, $f = \varinjlim f_i : X \longrightarrow Y$, Loc (8) tells us that if the $f_i \in W$, then $f \in W$. Here $I$ is
\msp{page 72}
an arbitrary (small) filtering category, and after replacing it by a cofinal category if necessary, we may suppose $I$ an ordered set. The axiom Loc (8 bis) corresponds to the case where $I = \mathbf{N}$.

Note that thanks to the construction of the category of paths, we have a factorization, functorial in $f$, of every arrow $f : X \longrightarrow Y$ of Cat as $f = pi$:
\[
X \xrightarrow{\quad i = i_f\quad} \bar{X}(f) = \bar{X} \xrightarrow{\quad p = p_f\quad} Y,
\]
with $i \in W_\omega$, $p \in \Parf_{W_\omega}$, a fortiori $i \in W$, $p \in \Parf_W$. This functor $f \longmapsto (\bar{X}(f), p_f)$ commutes with direct limits. Thus, if we apply this to each of the $f_\alpha : X_\alpha \longrightarrow Y_\alpha$, factoring it as
\[
X_\alpha \xrightarrow{i_\alpha} \bar{X}_\alpha \xrightarrow{p_\alpha} Y_\alpha,
\]
we find a direct system of diagrams of type $\boldsymbol{\Delta}^2$ (factoring the $f_i$), and on passing to the $\varinjlim$ we find
\[
X \xrightarrow{i} \bar{X} \xrightarrow{p} Y,
\]
which is none other than the canonical factorization
\msp{page 73}
of $f = \varinjlim f_i$ as $pi$, with $i \in W_\omega$, $p \in \Parf_{W_\omega}$ (N.B. Whereas neither $W_\omega$, nor $\Parf_{W_\omega}$ are stable under filtering $\varinjlim$, for sure\,!). We therefore have $i \in W$, and the condition $f \in W$, \textit{i.e.} $pi \in W$, is therefore equivalent to $p \in W$. This reduces us, for checking the axiom Loc (8) (resp. Loc (8 bis)), to the case where the $f_i$ are in $\Parf_{W_\omega}$.

Suppose now that $W$ satisfies $\mathrm{D}_{W_\omega}$ (the weakest of all the variants considered here of the axiom of $W$-fibrations Loc (7), with the sole exception of F). Then we have
\[
p \in W \qquad\quad \Longleftrightarrow \qquad\quad \text{the fibers of $p$ are $W$-aspheric,}
\]
\printedpage{41}
% head: A. Grothendieck : Dérivateurs, Ch. XVI, Édition des Universités de Montpellier II et Paris VII
the implication $\Longrightarrow$ resulting from the hypothesis $\mathrm{D}_{W_\omega}$ $\Parf_{W_\omega} \cap W \subset W^{\mathrm{u}}$, and $\Longleftarrow$ resulting from the sorites of $W$-smooth morphisms \textit{etc.}, via Loc (3). Likewise, the hypothesis that the $f_i \in W$ is equivalent to the one that the fibers of the $p_i$ are $W$-aspheric. On the other hand, it is immediate that
\msp{page 74}
every fiber of $p$ is isomorphic to the $\varinjlim\limits_{I}$ of the corresponding fibers of the $p_\alpha$. In this way, we find:

\textbf{Proposition 5.1.} \textit{Suppose that $W$ satisfies the axiom of $W$-fibrations} Loc (7 $\mathrm{D}_{W_\omega}$) (\textit{which is therefore the case if it satisfies one of the stronger axioms} Loc (7 a \textit{or} b \textit{or} c)). \textit{For $W$ to satisfy} Loc (8) (resp. \textit{to satisfy} Loc (8 bis)), \textit{it is necessary and sufficient that the set of objects in} Cat \textit{which are $W$-aspheric be stable under filtering $\varinjlim$} (resp. \textit{under $\varinjlim$ of countable sequences}).

5.2. Note that the category Cat is a paratopos, a fortiori it is \textit{accessible} (\textit{cf.} XVIII, §8), and the same therefore holds for $\uFl(\mathrm{Cat})$. Thus we have in $\mathrm{Fl}(\mathrm{Cat}) = \mathrm{Ob}\,\uFl(\mathrm{Cat})$ the notion of an \textit{accessible subset.} We shall therefore say that a subset $W$ of $\mathrm{Fl}\,\mathrm{Cat}$ (more generally of a category $\mathrm{Fl}(\mathcal{M})$, where $\mathcal{M}$ is an accessible category) is \textit{accessible,} if it is so with respect to the ambient category
\msp{page 75}
$\uFl\,\mathrm{Cat}$. This being so, the axiom Loc (9) is the following.

\textbf{Loc (9)}\quad $W$ is an \textit{accessible} subset of $\mathrm{Fl}(\mathrm{Cat}) = \mathrm{Ob}\,\uFl(\mathrm{Cat})$.

5.3. Let us recall that this somewhat technical notion simplifies considerably, if we suppose that $W$ satisfies the limits axiom Loc (8). Then $W$ is accessible if and only if there exists a \textit{small} subset $W_0$ of Fl(Cat), such that $W$ is the ``L$_1$-envelope'' of $W_0$, \textit{i.e.} the smallest subset of Fl(Cat) containing $W_0$ and L$_1$-stable (\textit{i.e.} stable under filtering $\varinjlim$). If this is so, it is clear that we can take $W_0$ of the form
\[
W_\pi = W \cap \mathrm{Fl}_\pi(\mathrm{Cat}),
\]
where $\pi$ is a suitable infinite cardinal, and where $\mathrm{Fl}_\pi(\mathrm{Cat}) \overset{\text{def}}{=} \mathrm{Fl}(\uCat_\pi)$, with $\uCat_\pi$ denoting the strictly full subcategory of Cat consisting of the objects $X$ such that $\mathrm{card}\,X \leq \pi$ (where by definition, $\mathrm{card}\,X = \mathrm{card}\,\mathrm{Fl}\,X$).

\msp{page 76}
Note that the filtration we have just described on Cat (indexed by the small infinite cardinals) is none other than its canonical filtration, \textit{i.e.} $\uCat_\pi$ is the full subcategory of Cat consisting of the $\pi$-accessible objects of Cat, and likewise in $\uFl(\mathrm{Cat})$. We thus find, thanks to the developments in XVIII, §13, the following characterization, which I recall here for reference (restricting myself for simplicity to the case where $W$ satisfies Loc (8)):

\textbf{Proposition 5.3.1.}\footnote[38]{\textit{cf. loc. cit.} 13.36.} \textit{Let $W$ be a subset of} Fl(Cat) \textit{stable under filtering $\varinjlim$, and let, for every infinite cardinal $\pi$, $W_\pi = W \cap \mathrm{Fl}_\pi(\mathrm{Cat})$, and $\underline{W}$, $\underline{W}_\pi$ be the corresponding full subcategories of} $\uFl(\mathrm{Cat})$. \textit{We say that $W$ is} $\pi$-accessible \textit{if $W$ is the L-envelope of $W_\pi$ (thus $W$ is accessible if and only if there exists an infinite cardinal $\pi$ such that $W$ is $\pi$-accessible). This being so, the following conditions are equivalent:}

\printedpage{42}
% head: A. Grothendieck : Dérivateurs, Ch. XVI, Édition des Universités de Montpellier II et Paris VII

\begin{enumerate}[label={}]
\item[a)] \textit{$W$ is $\pi$-accessible.}
\item[b)] \textit{For every direct system $(f_i)_{i\in I}$,}
\msp{page 77}
\textit{$L_1$-$\pi$-adapted in} $\uFl\,\mathrm{Cat}$ (cf. \textit{below}), \textit{with direct limit $f$, we have $f \in W$ if and only if the full subcategory $I_W$ of the $i \in \mathit{Ob}I$ such that $f_i \in W$ is cofinal in $I$.}
\item[b$'$)] (\textit{When we have chosen, for every $f \in$} Fl(Cat)\textit{, an $L_1$-$\pi$-adapted representation of $f$ as a $\varinjlim$}): \textit{If $f \in W$, then the full subcategory $I_W$ defined in} b) \textit{is cofinal in $I$.}
\end{enumerate}

Let us recall here the notion of an L$_1$-$\pi$-adapted direct system. The conditions are:
\begin{enumerate}[label={}]
\item[1°)] $I$ is large relative to $\pi$, and stable under filtering $\varinjlim$ of cardinal $\leq \pi$.
\item[2°)] The functor $i \longmapsto x_i : I \longrightarrow \uFl(\mathrm{Cat})$ commutes with the $\varinjlim$ in question.
\item[3°)] The $x_i$ are in $\uFl(\mathrm{Cat})_\pi \overset{\text{def}}{=} \uFl_\pi(\mathrm{Cat})$.
\end{enumerate}
Note that in the case of a ``parabelian paratopos'' such as Cat or $\uFl\,\mathrm{Cat}$, there is a canonical way, $\pi$ being given, to represent an object $f$ as the $\varinjlim$ of an L$_1$-$\pi$-adapted direct system. Here,
\[
f : X \longrightarrow Y
\]
being given, we consider the set $I$ of
\msp{page 78}
pairs $(X_i \subset X, Y_i \subset Y)$ of subcategories of $X$ and $Y$, such that we have $X_i, Y_i \in \mathrm{Ob}\,\mathrm{Cat}_\pi$ (\textit{i.e.} $\mathrm{card}\,(X_i), \mathrm{card}\,(Y_i) \leq \pi$) and that $f(X_i) \subset Y_i$. It is immediate that $I$ is filtering, for the obvious order relation
\[
(X_i, Y_i) \subset (X_j, Y_j) \qquad\quad \Longleftrightarrow \qquad\quad X_i \subset Y_i \ \text{ and } \ X_j \subset Y_j, % [sic] X_i \subset Y_i et X_j \subset Y_j -> X_i \subset X_j et Y_i \subset Y_j, la relation d'ordre compare les couples (X_i,Y_i) et (X_j,Y_j)
\]
and even large relative to $\pi$, that $I$ is stable under filtering $\varinjlim$ of cardinal $\leq \pi$, and that the functor
\[
i \longmapsto (f_i : X_i \longrightarrow Y_i) : I \longrightarrow \uFl\,\mathrm{Cat}
\]
commutes with the said $\varinjlim$ , finally the $f_i$ are in $\uFl(\mathrm{Cat})_\pi$ by construction. Thus we have an L$_1$-$\pi$-adapted direct system, and it is clear that its direct limit is canonically isomorphic to $f$. This having been seen:

\textbf{Corollary 5.3.2.} \textit{The condition that $W$ be $\pi$-accessible therefore means that for every $f : X \longrightarrow Y \in W$, and for every subcategory $X_i$ of $X$ and $Y_i$ of $Y$, of cardinal $\leq \pi$ (with $f(X_i) \subset Y_i$), but in fact this is unnecessary, after enlarging $Y_i$ if necessary), there exist subcategories} % [sic] (avec f(X_i) \subset Y_i), mais ... agrandir Y_i) -> (avec f(X_i) \subset Y_i, mais ... agrandir Y_i), parenthèse fermante surnuméraire
\msp{page 79}
\textit{$X_j \supset X_i$ and $Y_j \supset Y_i$ of cardinal $\leq \pi$, such that $f$ induces $f_j : X_j \longrightarrow Y_j$ and that we have $f_j \in W$.}

The notion of accessibility of $W$ seems to me as well understood as that of limits, I no longer see any mystery in it. On the other hand, it seems to me that all the basic localizers used to date satisfy it (just as they do the axiom Loc (8)). Here, I shall confine myself to establishing:

\printedpage{43}
% head: A. Grothendieck : Dérivateurs, Ch. XVI, Édition des Universités de Montpellier II et Paris VII

\textbf{Theorem 5.4.}\footnote[39]{Well deserves the name of \textit{theorem,} especially if we include in it the whole list of Loc axioms mentioned opposite\,! This theorem is not at all obvious, starting from the cohomological definition of $\mathbf{W}_\infty$ - it is a result (``well known'') which seems deep to me.} \textit{The basic localizer $\mathbf{W}_\infty$ (``usual'' weak equivalences) satisfies the axiom} Loc (9) (\textit{in fact, it satisfies} all \textit{the axioms seen so far, namely} Loc (1 ter, 2, 3 bis, 4, 5, 6, 7, 8, 9)). \textit{More precisely, $\mathbf{W}_\infty$ is $\aleph_0$-accessible (where $\aleph_0$, or $\pi_0$, is the first infinite cardinal).} % [sic] les axiomes vu -> les axiomes vus, accord du participe

I do not return here to the verification of the other axioms, and confine myself to proving Loc (9), accessibility, using what may be useful to me. I
\msp{page 80}
confine myself to indicating the proof (rather technical) in its broad outlines. For the moment, $W$ is again an arbitrary basic localizer, except that it satisfies Loc (8).

\textbf{Lemma 5.5.}\footnote[40]{We suppose that $W$ satisfies Loc (8).}
\begin{enumerate}[label={}]
\item[a)] \textit{For the criterion} 5.3.1 b) \textit{of $\pi$-accessibility of $W$ to be satisfied, it is necessary and sufficient that it be so when the $f_i$ are $W_\omega$-perfect.}
\item[b)] \textit{Suppose that $\mathit{Parf}_{W_\omega} \cap W \subset W^u$,} i.e. \textit{that $W$ satisfies} Loc (7 $\mathrm{D}_{W_\omega}$) (\textit{the weakest of the variants of the axiom of $W$-fibrations, except only for} Loc (7 F)). \textit{Suppose moreover (to be on the safe side) that $W$ satisfies} \texttt{[unfinished sentence]}. \texttt{[?Then]} \textit{$W$ is $\pi$-accessible if and only if the set} $\mathrm{Asph}_W \subset \mathrm{Ob}\,\mathrm{Cat}$ \textit{of $W$-aspheric objects of} Cat \textit{is $\pi$-accessible. Or again, that the subset} $\mathrm{Asph}_W \subset \mathrm{Ob}\,\mathrm{Cat}_{\textit{0-connexe}}$ \textit{be $\pi$-accessible, where} $\mathrm{Cat}_{\textit{0-connexe}}$ \textit{denotes the full subcategory of} Cat \textit{consisting of} 0-\textit{connected objects.}
\end{enumerate}
Essentially, the proof is the same as that of 5.1, a little more technical. The essential tool is still the functorial factorization $f = pi$ of an arrow of Cat, with $i \in W_\omega$, $p \in \Parf_{W_\omega}$, and the property of
\msp{page 81}
commutation with filtering $\varinjlim$ of the said factorization.

Let us now return to the case of $W = \mathbf{W}_\infty$. We must prove that $\mathrm{Asph}_{\mathbf{W}_\infty} \subset \mathrm{Ob}\,\mathrm{Cat}_{\text{0-connected}}$ is $\pi_0$-accessible. Thus to prove the criterion corresponding to 5.3.1 b), placing ourselves now in $\mathrm{Cat}_{\text{0-connected}}$ instead of $\mathrm{Fl}\,\mathrm{Cat}$. We therefore have
\[
X = \varinjlim_{i\in I} X_i,
\]
an L$_1$-$\pi$-admissible representation of $X$ $\mathbf{W}_\infty$-aspheric, to prove that for every $i$ in $I$, $\exists j$ in $I$ and $i \longrightarrow j$, with $X_j$ $\mathbf{W}_\infty$-aspheric.

I admit here the AQT \texttt{[âne qui trotte: at a donkey's trot, i.e. by routine verification]} lemma:

\textbf{Lemma 5.6.} \textit{Let $\mathbf{W}_0 = \{f \in \mathrm{Fl}\,\mathrm{Cat} \mid \pi_0(f) \ \ \textit{isomorphisme}\}$. Then $\mathbf{W}_0$ is $\pi_0$-accessible (and $L_1$-stable). A fortiori, the set of} 0-\textit{connected objects of} Cat \textit{is $\pi_0$-accessible.}

This reduces us, after replacing $I$ by a cofinal category if necessary, to supposing that the $X_i$ are 0-connected. (But I forgot that this reduction has already been made\,! So one must view 5.6 as an ingredient